\chapter{文件系统、程序装载与 shell}

\section{最简内存文件系统格式}
文件系统使用扁平名称空间、固定数量文件项和固定大小数据槽。镜像由文件系统头、固定长度文件项表、对齐填充和数据槽组成，所有多字节字段使用小端编码。

\section{文件操作语义}
学生实现创建、删除、整文件读取和整文件覆盖写入。名称重复、文件不存在、表满、内容超限、缓冲不足或非法指针失败时不得产生可观察的部分修改。

\section{最小 ELF 与按名称启动}
仅支持静态 RV64 \code{ET\_EXEC} ELF 的 \code{PT\_LOAD} 段。shell 通过文件名启动前台程序；完成事件经专用消息队列通知，不提供 fork、waitpid、管道或作业控制。

\section{不纳入本版的能力}
目录、路径、fd 文件接口、VFS、动态链接、按需分页、swap、mmap/munmap、管道、重定向、后台任务与完整 shell 语法均不属于基础必做范围。
